\section{Setting, notation and data}\label{sec:setting}

\subsection{The decision problem}

An operator scores a stream of events $x_1,x_2,\dots$ arriving in temporal
order. A model $f$ maps an event to $p(x)\in[0,1]$, and a threshold $\tau$ maps
the score to an action $a\in\{0,1\}$. Errors are asymmetric: a false negative
costs $c_{\mathrm{FN}}$, a false positive costs $c_{\mathrm{FP}}$. The quantity
we optimise is the expected cost per event,
\begin{equation}
C(\tau) \;=\; c_{\mathrm{FN}}\,\pi\,\big(1-\mathrm{TPR}(\tau)\big)
\;+\; c_{\mathrm{FP}}\,(1-\pi)\,\mathrm{FPR}(\tau),
\label{eq:cost}
\end{equation}
with $\pi$ the base rate. Two consequences shape every experiment below.
First, the operating point matters more than the ranking: a model with a better
AUC-PR can be worse in \eqref{eq:cost} at the deployed threshold. Second, the
threshold must be chosen without looking at the test period, so we choose $\tau$
once on a validation split for each cost ratio and evaluate once on test.
Throughout, AUC-PR is the primary ranking metric because the positive class is
rare ($\pi\approx\BafFraudPct\%$ on \DatasetBaf), AUC-ROC is secondary, and we
report Recall at a fixed $1\%$ false-positive rate as the auditable operating
point, following \cite{elkan2001foundations}.

\subsection{The mask family}\label{sec:family}

Let $i$ index queries and $j$ keys, both in $\{0,\dots,n-1\}$. The family we
analyse is
\begin{equation}
\Sset(i) \;=\; \underbrace{\{j : i-w < j \le i\}}_{\text{window}}
\;\cup\; \underbrace{\{j : j \le i,\ j \bmod g = 0\}}_{\text{periodic anchors}}
\;\cup\; \underbrace{\{j : j < s\}}_{\text{sinks}},
\label{eq:mask}
\end{equation}
with every set intersected with the causal constraint $j\le i$. The three
parameters specialise to the designs in use:

\begin{itemize}
\item $g=\infty$, $s=0$: a pure sliding window. Global-key budget zero.
\item $g=\infty$, $s>0$: a window plus a fixed number of sink positions; this
is the mechanism by which a fixed set of global tokens is usually implemented.
\item $g$ fixed and finite: periodic anchoring, i.e. every $g$-th position is
visible to all later queries. This is the design we analyse and measure; it is
the natural way to guarantee that no position is more than $g$ steps from a
globally visible key, and it appears whenever a system wants coverage rather
than a hand-chosen global set.
\item $w=n$, $g=\infty$: dense causal attention, the reference.
\end{itemize}

The distinction between the second and third items is the subject of
Theorem~\ref{thm:linear}: both are described as ``global attention'', but the
second has a constant number of global keys and the third has $n/g$ of them.

\subsection{Data}

We use two public benchmarks, both with source-provided temporal splits, and we
never shuffle across time.

\DatasetIeee{} \cite{ieeecis2019} contains \IeeeRows{} card transactions with
\IeeeCols{} columns and a fraud rate of \IeeeFraudPct\%. The source provides
strict temporal partitions of \TrainRows{}, \ValRows{} and \TestRows{} rows. The
mirror we use is checksummed against the published dataset-card hash (first
sixteen hex digits \IeeeSha) before any experiment runs.

\DatasetBaf{} \cite{jesus2022turning} contains \BafRows{} account-opening
records with \BafCols{} columns and a rate of \BafFraudPct\%. It is
source-declared synthetic, which we treat as a limitation rather than a
convenience: it is the second dataset, never the primary one.

\paragraph{Entity sequences} A decision is often better made about a card than
about a single transaction. Grouping the mirror by \texttt{card1} and ordering by
timestamp yields \NEtities{} entities whose length distribution is heavy-tailed:
median \MedianLen, 90th percentile \PninetyLen, 99th percentile
\PninetyNineLen, maximum \MaxLen, with \EntGteThousand{} entities reaching
$n\ge 1024$. This distribution is itself a result and it constrains the
attention question: long-context machinery matters for a thin tail of entities,
not for the median one. We therefore report the attention measurements on
synthetic sequences of controlled length and use the real sequences only for the
autocorrelation analysis of Section~\ref{sec:acf} and the retained-mass check of
Section~\ref{sec:mass}.

\paragraph{Feature budget} To keep the comparison honest across model families
we fix a budget of \MainFeatures{} numeric predictors, \MainTrainRows{} training
rows and \MainTestRows{} held-out rows, with the operating point chosen on
validation. Limiting the budget below the data's capacity is deliberate: it
makes the supervised baselines weaker than they could be, which biases the
comparison \emph{in favour} of the pretrained model, and is therefore the
conservative direction for claim C5.

\subsection{Hardware and software}

All device measurements come from one \GpuName{} with \GpuMemGiB\,GiB of memory
and compute capability \GpuCapability, using PyTorch \TorchVer{} built against
CUDA \CudaVer, float16 activations, block size $B=\GpuBlock$, window
$w=\GpuW$ and anchor period $g=\GpuG$, with \GpuReps{} timed repetitions per
cell after warm-up. Tabular models and the statistical layer run on the host
CPU; the on-chain cost measurement runs on a local EVM (\SolcVer) with no
network. We report single-device numbers because the deployment target is a
single device, and we say so in every caption where it matters.
